\section{Runtime Co-Design from the Measured Bottlenecks}
\label{sec:runtime}
\pending\ This section reserves approximately 1.5--2 pages of the ten-page planning budget for a possible future measured contribution. It contains requirements, hypotheses, an interface proposal, and evaluation questions only. The hardware study is complete without it. No speed, energy, correctness, implementation, or ablation result is attributed to the proposed design.

\subsection{Requirements traced to closed measurements}
The direct stateful measurements motivate a narrow intervention: retain usable draft state across speculative rounds and reconcile it with the committed target history. The NPU replay-plus-proposal region motivates this requirement; the tiny measured residual does not motivate a broad rewrite of the orchestrator. The CPU result additionally warns that reducing draft cost alone may leave complete target-side work dominant. State persistence is therefore a hypothesis about a measured component, not a forecast of overall profitability.

The verification controls impose a separate requirement. A runtime must qualify the target path under its chosen output contract before using that path's timings to select an execution policy. A cheaper blocked path is not a candidate optimization under the current contract. B570 further requires that serialization be recorded as an execution setting, with exactness independently checked across the intended workload. A small probe cannot certify an unrestricted path selection policy.

The M4 contrast motivates an optional profitability guard based on realized quantities. Its purpose would be to decide whether a qualified speculative path is useful for the current workload, using direct committed progress and complete costs. It must not infer realized draft time from an isolated throughput number or count accepted draft tokens as all output progress. The proposed minimum study holds proposal width and device placement fixed; adaptive width and heterogeneous device selection would add separate questions only if later evidence justifies them.

\subsection{Minimal state interface proposal}
We propose a request-scoped draft-state object associated with a committed token prefix and a model/runtime configuration. A begin operation would initialize state from the prompt. A checkpoint operation would identify the committed boundary before proposing. A propose operation would return candidate token IDs and a state handle describing their relation to that boundary. A reconcile operation would retain only the accepted draft prefix, incorporate any target correction or bonus token, and discard rejected candidate state. These are proposed interface responsibilities, not available functions in the measured adapter.

The target reference and its correctness contract would stay explicit. After every round, the target logical history, draft logical history, and committed output ledger would have to agree. Candidate state cannot become committed state merely because a proposal call completed. At request termination, all handles must be invalidated so a later request cannot inherit another prefix. If the backend cannot expose safe truncation or restoration, a replay path may be necessary; any such fallback cost must stay inside the realized draft boundary.

The interface should expose a complete draft-call timer, complete target-round timer, residual interval, and committed-token count under one enclosing decode window. Checkpoint creation, state copying, cache allocation, synchronization, and recovery must be charged to the responsible measured region. Merely moving work outside a narrow proposal timer would not establish an improvement. If replay and proposal remain inseparable, the evaluation should keep the aggregate region used by the hardware study rather than claiming an unavailable subcomponent decomposition.

The proposal does not require a new speculative algorithm, a new kernel, or a new learned draft. Its central test is whether the qualified backend can preserve useful state without changing token semantics or paying comparable reconciliation costs. Memory capacity and state-handle lifetime are therefore evaluation questions alongside latency. Persistent state and KV reuse have substantial prior art; their existence is not the intended novelty claim.

\subsection{Correctness-qualified path selection}
A minimal path registry would associate each target execution setting with its qualification scope: model artifacts, runtime revision, numerical options, context/output conditions, and the evaluated workload. Selection would consider only paths admitted by the active contract. Changing the runtime, quantization, or execution shape would invalidate the relevant qualification rather than inherit it silently. This registry is proposed metadata and policy, not a claim that a finite test suite proves correctness for every possible input.

The evaluation must retain mismatches and blocked candidates. A path that fails remains a diagnostic observation with its failure scope visible. Recovery to target-only execution must preserve the same committed history and have its cost included. A future study could consider a different approximate-quality contract, but it would need its own quality metric and protocol; it cannot retrospectively relabel the frozen failures in this manuscript.

\subsection{Planned evaluation questions and decision rule}
A separate protocol should be frozen before observing intervention results. It should compare the current replay adapter, the proposed persistent-state adapter, and matched target-only decoding on the same prospective workload and timing boundary. Correctness qualification should exercise accepted prefixes, rejection, correction, bonus output, request reset, and output-budget boundaries. Full-workload exactness must accompany any promoted speed result; passing an initial gate alone is insufficient.

The key questions are whether the interface preserves token equivalence; whether realized draft cost decreases after all state-management work is counted; whether complete target-side cost changes; and whether the end-to-end prompt-level comparison improves. Component measurements and total-time accounting should reconcile. An energy claim would require a newly qualified end-to-end sensor boundary, which the closed 358H campaign does not supply.

If state persistence reduces draft cost but leaves the system unprofitable, the result should be reported in that form. If correctness fails, the timing remains diagnostic. If the implementation is unavailable or yields no defensible contribution, this entire section can be removed without changing the hardware-study conclusion. A successful future intervention would need measured evidence and a revised abstract, contribution list, and conclusion; none is assumed in manuscript v0.1.
